% Source: https://artofproblemsolving.com/community/c3h2309215_100_posts_problems
% Stable identity: aops-topic-2309215-q08
Let $O$ be the centre and $A_0,\dots,A_4$ the vertices, so $OA_k=2$. You may use $\cos36^\circ=\dfrac{1+\sqrt5}{4}$ and $\cos72^\circ=\dfrac{\sqrt5-1}{4}$.

\textbf{The inner vertices.} Let $Q$ be the crossing point of the diagonals $A_0A_2$ and $A_1A_4$, the inner vertex of the star between the points $A_0$ and $A_1$. The reflection in the perpendicular bisector of $A_0A_1$ swaps $A_0\leftrightarrow A_1$ and $A_2\leftrightarrow A_4$, so it swaps the two diagonals and fixes $Q$: hence $Q$ lies on the ray from $O$ through the midpoint of $A_0A_1$, which makes an angle of $36^\circ$ with $OA_0$. The chord $A_1A_4$ subtends the angle $144^\circ$ at $O$, so its distance from $O$ is $2\cos72^\circ$, measured along $OA_0$ (which is perpendicular to $A_1A_4$ by symmetry). Since $OQ$ makes the angle $36^\circ$ with this perpendicular,
\[OQ\cos36^\circ=2\cos72^\circ,\qquad OQ=\dfrac{2\cos72^\circ}{\cos36^\circ}=\dfrac{2(\sqrt5-1)}{\sqrt5+1}=\dfrac{2(\sqrt5-1)^2}{4}=3-\sqrt5.\]

\begin{asy}
usepackage("amsmath");
size(6cm);
defaultpen(fontsize(10pt));
real R=2, rho=3-sqrt(5);
pair[] P; pair[] Q;
for(int k=0;k<5;++k){ P.push(R*dir(90+72k)); Q.push(rho*dir(126+72k)); }
path st=P[0]--Q[0]--P[1]--Q[1]--P[2]--Q[2]--P[3]--Q[3]--P[4]--Q[4]--cycle;
fill(st,rgb(0.85,0.92,1));
draw(circle((0,0),R),gray(0.6));
draw(P[0]--P[1]--P[2]--P[3]--P[4]--cycle,gray(0.4));
draw(P[0]--P[2]--P[4]--P[1]--P[3]--cycle,linewidth(0.9));
pair O=(0,0);
fill(O--P[0]--Q[0]--cycle,rgb(1,0.8,0.7));
draw(O--P[0],red); draw(O--Q[0],red); draw(P[0]--Q[0],red+linewidth(0.9));
draw(O--0.95*R*dir(126),dashed+gray(0.4));
dot(O); label("$O$",O,S);
for(int k=0;k<5;++k) dot(P[k]);
label("$A_0$",P[0],N); label("$A_1$",P[1],W); label("$A_2$",P[2],SW); label("$A_3$",P[3],SE); label("$A_4$",P[4],E);
dot(Q[0],red); label("$Q$",Q[0],dir(222),red);
label("$2$",(O+P[0])/2,E,red);
draw(arc(O,0.28,90,126),red);
label("$\angle A_0OQ=36^\circ$",(0,-2.05),S,fontsize(9pt));
\end{asy}

\textbf{The area.} By symmetry the star consists of $10$ triangles congruent to $OA_0Q$, each with sides $OA_0=2$, $OQ=3-\sqrt5$ and the angle $36^\circ$ between them. Also
\[\sin36^\circ=\sqrt{1-\left(\dfrac{1+\sqrt5}{4}\right)^2}=\sqrt{\dfrac{10-2\sqrt5}{16}}=\dfrac{\sqrt{10-2\sqrt5}}{4}.\]
So the area is
\[10\cdot\dfrac12\cdot2\,(3-\sqrt5)\sin36^\circ=\dfrac52(3-\sqrt5)\sqrt{10-2\sqrt5}.\]
Since $(3-\sqrt5)^2(10-2\sqrt5)=(14-6\sqrt5)(10-2\sqrt5)=200-88\sqrt5$, this equals
\[\dfrac52\sqrt{200-88\sqrt5}=\boxed{5\sqrt{50-22\sqrt5}}\approx4.49.\]
